\section{Project context and objectives}

\subsection{Background}
Running applications are effective at measuring distance and pace, while travel
applications are effective at describing places. These experiences are usually
separate: a run becomes a graph with little cultural meaning, and a destination
guide rarely turns exploration into an active, personal challenge. Questory
connects the two. Movement provides the structure, landmarks provide motivation,
and photographs become evidence and material for a visual diary.

The product was designed for an academic mobile-development project, but its
scope is intentionally closer to a complete small application than a screen
prototype. It has persistent local data, real location and camera integrations,
error and permission states, navigation across the full journey, deterministic
content, an editable export workflow, and an optional online extension.

\subsection{Target users}
\begin{itemize}
  \item recreational runners who want a motivating route without performance
        pressure or a public leaderboard;
  \item visitors who prefer active exploration and need useful content even when
        mobile data is weak or unavailable;
  \item local residents who want a creative record of familiar places;
  \item privacy-conscious users who do not want an account for basic use; and
  \item student evaluators who need a repeatable, inspectable end-to-end flow.
\end{itemize}

\subsection{Problem statement}
The core problem is to record a credible run while helping the user notice the
city and leave with something more expressive than a route screenshot. The
solution must tolerate permission denial, disabled GPS, noisy points, temporary
camera files, process interruption, missing network access, and image-export
failure without losing the user's local activity.

\subsection{Objectives and design principles}
Questory's objectives are to provide useful bundled routes for two Vietnamese
cities, support both structured and self-directed running, evaluate proximity
quests fairly, preserve unfinished and completed work locally, and export a
shareable portrait artifact. Five principles shaped the implementation:

\begin{enumerate}
  \item \textbf{Offline is the default.} Network access cannot be a prerequisite
        for opening the app, recording, finishing, reviewing, or exporting.
  \item \textbf{User action controls sensors.} Tracking begins only after an
        explicit explanation and confirmation.
  \item \textbf{State is recoverable.} Active runs are checkpointed and retained
        photos move into application-controlled storage.
  \item \textbf{The output remains editable.} Story Studio saves a serializable
        document rather than only a rendered image.
  \item \textbf{Optional systems remain optional.} Remote sharing is behind an
        interface and disappears gracefully when not configured.
\end{enumerate}

\subsection{Originality and scope}
Questory's distinctive element is the complete loop from place discovery to
physical activity, evidence capture, personal achievement, and editable visual
story. It is not a map viewer with decorative screens and not a camera gallery
with a running label. The route, GPS track, quest evidence, summary, and story
document share stable identities and survive navigation and application restart.

